% Compile from this directory with:
% latexmk -pdf -interaction=nonstopmode -halt-on-error sr_vector_magnetometry_model.tex
\documentclass[11pt,letterpaper]{article}
\usepackage[margin=0.85in,headheight=14pt]{geometry}
\usepackage[T1]{fontenc}
\usepackage{lmodern}
\usepackage{amsmath,amssymb,bm}
\usepackage{booktabs,tabularx,array,microtype}
\usepackage{enumitem,fancyhdr}
\usepackage[hidelinks]{hyperref}
\newcommand{\Tr}{\operatorname{Tr}}
\newcommand{\id}{\mathbb{I}}
\newcommand{\OD}{\mathrm{OD}}
\newcommand{\dt}{\delta t}
\newcommand{\Je}{J_{\epsilon}}
\newcommand{\Jk}{J_k}
\newcommand{\Sr}{${}^{88}\mathrm{Sr}$}
\newcommand{\Ptwo}{${}^{3}P_2$}
\newcommand{\Done}{${}^{3}D_1$}
\newcolumntype{Y}{>{\raggedright\arraybackslash}X}
\setlength{\parindent}{0pt}
\setlength{\parskip}{5pt}
\setlist[itemize]{leftmargin=1.4em,itemsep=3pt,topsep=3pt}
\renewcommand{\arraystretch}{1.2}
\pagestyle{fancy}
\fancyhf{}
\fancyhead[L]{\small Sr vector magnetometry}
\fancyhead[R]{\small Model discussion, 27 September 2026}
\fancyfoot[C]{\thepage}
\begin{document}

\begin{center}
{\LARGE\bfseries Vector magnetometry with metastable strontium}\\[7pt]
{\large Tensor control, probe-induced loss, and continuous readout}\\[6pt]
{\small Discussion record based on the handwritten note of 26 September 2026}
\end{center}

\textbf{Status.} This document records the model discussion and its consistency
checks. It proposes a starting model; it does not report a new simulation or an
agreed implementation. The existing v7 simulation is unchanged.

\section{Physical proposal and atomic data}

Prepare an initially uncorrelated ensemble of \Sr\ atoms in the metastable
\Ptwo\ manifold, with electronic angular momentum $J=2$ and five magnetic
sublevels. An unknown magnetic field drives Larmor precession. A known magnetic
control and the tensor light shift of a linearly polarized probe modify the
spin dynamics. Faraday polarimetry of that same probe produces a continuous
measurement record~\cite{note}.

The probe is detuned from the ${}^{3}P_2\rightarrow{}^{3}D_1$ transition.
Figure~1(c) of Ref.~\cite{masson} shows the relevant levels and decay paths;
Table~II gives the following single-atom rates for \Sr.

\begin{center}
\begin{tabular}{@{}lrr@{}}
\toprule
Decay channel & Wavelength (nm) & Partial rate ($\mathrm{s}^{-1}$)\\
\midrule
${}^{3}D_1\rightarrow{}^{3}P_0$ & 2600 & $2.8\times10^5$\\
${}^{3}D_1\rightarrow{}^{3}P_1$ & 2740 & $1.8\times10^5$\\
${}^{3}D_1\rightarrow{}^{3}P_2$ & 3070 & $8.8\times10^3$\\
\bottomrule
\end{tabular}
\end{center}

Summing these channels gives
\begin{align}
\Gamma &=4.688\times10^5\ \mathrm{s}^{-1},
& \frac{\Gamma}{2\pi}&\simeq74.6\ \mathrm{kHz},\label{eq:linewidth}\\
b_{\mathrm{return}}&=\frac{8.8\times10^3}{\Gamma}\simeq0.01877,
& b_{\mathrm{loss}}&\simeq0.98123.\label{eq:branching}
\end{align}
Thus about $98.12\%$ of scattering events leave the sensing manifold and
$1.88\%$ return. This supports the note's loss-dominated approximation, with a
small correction to its rounded statement of $99\%$ loss. ``Loss'' means loss
from \Ptwo, not necessarily physical escape from the trap. Treating these atoms
as dark also assumes negligible repumping and negligible response of the other
manifolds at the chosen probe frequency.

\textbf{Scope of the paper comparison.} We use the atomic structure and
single-atom rates from Ref.~\cite{masson}. The magnetometry model considered here
assumes independent atoms, without the collective radiative interactions and
superradiance studied in that paper.

\newpage
\section{Coherent dynamics and parameter conventions}

Set $\hbar=1$, so Hamiltonians have units of angular frequency. Let
$\boldsymbol\epsilon$ be the real unit polarization vector of the probe and
$\bm k$ its unit propagation direction. Define
\begin{equation}
\Je=\boldsymbol\epsilon\cdot\bm J,
\qquad \Jk=\bm k\cdot\bm J,
\qquad \boldsymbol\epsilon\cdot\bm k=0.
\end{equation}
The tensor shift acts along the polarization axis; Faraday rotation measures
the spin along the propagation axis. For example, keeping the current code's
tensor term $J_z^2$ and readout $J_y$ corresponds to
$\boldsymbol\epsilon=\bm e_z$ and $\bm k=\bm e_y$. A readout of $J_z$ instead
requires a transverse polarization and the corresponding tensor axis.

Using the magnetic-field sign convention of the handwritten note,
\begin{equation}
\boxed{
H(t)=-\bm\omega\cdot\bm J
+\Omega\bigl[\cos\phi(t)J_x+\sin\phi(t)J_y\bigr]
+\frac{\beta}{2J}\Je^2,
\qquad J=2.
}
\label{eq:hamiltonian}
\end{equation}
Here $\bm\omega$ is the unknown Larmor-frequency vector, $\Omega$ is the known
magnetic-control amplitude, and $\phi(t)$ is a prescribed control phase. The
current v7 script uses $+\bm w\cdot\bm J$, so the conventions agree upon setting
$\bm w=-\bm\omega$. This sign choice must be consistent in data generation,
fitting, and conversion back to a magnetic field.

For an isolated, weakly excited optical transition, define
\begin{equation}
\Delta_L=\omega_L-\omega_{eg},\qquad
V_0=\frac{\Omega_L^2}{4\Delta_L},\qquad
\gamma_0=\frac{\Omega_L^2\Gamma}{4\Delta_L^2}.
\label{eq:scales}
\end{equation}
$\Omega_L$ is a reference optical Rabi frequency, distinct from $\Omega$.
For the angular factors below, it is normalized so that the actual absorption
Rabi frequency equals $\Omega_L$ times the corresponding Clebsch--Gordan
coefficient. All detunings and Rabi frequencies use angular-frequency units;
$\Gamma$ is the population decay rate in Eq.~\eqref{eq:linewidth}.

The relation to the simulation's tensor parameter is
\begin{equation}
\boxed{
\frac{\beta}{2J}=C^{(2)}V_0,
\qquad
\frac{\beta/(2J)}{\gamma_0}=C^{(2)}\frac{\Delta_L}{\Gamma}.
}
\label{eq:beta}
\end{equation}
For $C^{(2)}=-1/10$ and $J=2$, this gives
$\beta=-\Omega_L^2/(10\Delta_L)$, or equivalently
$\beta/\gamma_0=-(2/5)\Delta_L/\Gamma$.
Here $\beta$ is signed. v7 instead sweeps a positive magnitude and applies the
tensor sign separately through its atomic preset.

\textbf{Two normalization cautions.} The coefficient multiplying $\Je^2$ is
$\beta/(2J)$, not $\beta$. Also, $\gamma_0$ is a reference scattering-rate
scale, not the actual state-independent photon-scattering rate. The latter
requires the angular operator derived on the next page. The note's later
scaling argument suppresses these angular factors; they must be restored for
quantitative comparisons.

\newpage
\section{State-dependent loss and the tensor correction}

The handwritten coefficients are read as
\begin{equation}
C^{(0)}=\frac15,\qquad C^{(2)}=-\frac1{10}.
\end{equation}
In the irreducible-tensor convention, the rank-two operator includes a trace
subtraction~\cite{deutsch}. Therefore the absorption-strength operator is
\begin{align}
Q&=C^{(0)}\id+C^{(2)}
\left(\Je^2-\frac{J(J+1)}{3}\id\right)\nonumber\\
&=\frac15\id-\frac1{10}(\Je^2-2\id)
=\boxed{\frac{4\id-\Je^2}{10}}.
\label{eq:Q}
\end{align}
An independent angular-momentum check gives
$Q_{mm}=|\langle 2,m;1,0\mid1,m\rangle|^2$. Explicitly,
\begin{center}
\begin{tabular}{@{}lccccc@{}}
\toprule
$m_\epsilon$ & $-2$ & $-1$ & $0$ & $1$ & $2$\\
\midrule
$Q_{mm}$ & $0$ & $3/10$ & $2/5$ & $3/10$ & $0$\\
\bottomrule
\end{tabular}
\end{center}

The literal expression $\id/5-\Je^2/10$ in the note would give negative loss
rates for $m_\epsilon=\pm2$. The missing scalar contribution is essential.
One may discard a common \emph{real energy shift} from $H$, but cannot discard
a common decay term from its non-Hermitian part.

\subsection*{Starting approximation: every scattering event is loss}
Approximating the $98.12\%$ loss branch by unity gives
\begin{equation}
K=\gamma_0Q,\qquad
H_{\mathrm{eff}}=H-\frac{i}{2}K,\qquad
\boxed{\dot\rho=-i[H,\rho]-\frac12\{K,\rho\}.}
\label{eq:loss}
\end{equation}
The density matrix here is unnormalized, with $\Tr\rho(0)=1$:
\begin{equation}
p(t)=\Tr\rho(t),\qquad
\dot p(t)=-\gamma_0\Tr[Q\rho(t)].
\label{eq:survival}
\end{equation}
It describes the surviving part of the ensemble after ignoring the outgoing
states. Since $Q\ge0$, survival cannot increase. The instantaneous scattering
rate per surviving atom is $\gamma_0\Tr(Q\rho/p)$ when $p>0$.

\textbf{Physical consequence.} In the isolated-line approximation,
$m_\epsilon=\pm2$ are dark to the linearly polarized probe. Magnetic dynamics
can move population between dark and bright states, so survival generally is
not a single exponential independent of the unknown field and control.

\subsection*{Optional refinement: retain the small return branch}
An explicit return model would add $\sum_a L_a\rho L_a^\dagger$ to
Eq.~\eqref{eq:loss}, with
$\sum_a L_a^\dagger L_a=b_{\mathrm{return}}\gamma_0Q$.
The trace would then decrease at
$b_{\mathrm{loss}}\gamma_0\Tr(Q\rho)$. The $L_a$ must be derived from the
absorption and emission angular factors; v7's symmetric $J_\pm$ channels are
not automatically the correct return model.

\newpage
\section{Measurement record, survival, and shot noise}

Let $N_0$ be the initial atom number. Define the measurement strength $\kappa$
operationally: in a time bin $\dt$, the calibrated collective-spin readout has
shot-noise variance $1/(\kappa\dt)$. In the weak-backaction, shot-noise-dominated
approximation, the record per initial atom is
\begin{equation}
\boxed{
m_i=\Tr[\Jk\rho(t_i)]
+\frac{\xi_i}{N_0\sqrt{\kappa\dt}},
\qquad \xi_i\overset{\mathrm{iid}}{\sim}\mathcal N(0,1).
}
\label{eq:record}
\end{equation}
This treats each bin as sufficiently short to use its instantaneous signal;
otherwise the mean must be averaged across the bin. It models ensemble means
with additive photon shot noise, omitting resolved atom-number fluctuations,
projection noise, and measurement-induced collective correlations.

\subsection*{Survival appears exactly once}
Writing $\widetilde\rho=\rho/p$ for the normalized state of surviving atoms,
\begin{equation}
\Tr(\Jk\rho)=p(t)\Tr(\Jk\widetilde\rho).
\end{equation}
No additional survival factor belongs in Eq.~\eqref{eq:record}. Conversely,
if we replace the loss operator by $K=\gamma_{\mathrm{loss}}\id$, then
\begin{equation}
m_i=e^{-\gamma_{\mathrm{loss}}t_i}
\langle\Jk\rangle_{\mathrm{conditional},i}
+\frac{\xi_i}{N_0\sqrt{\kappa\dt}}.
\label{eq:uniform-record}
\end{equation}
Page~3 of the handwritten note includes the exponential in the collective
record $M(t_i)$ but omits it in $m(t_i)=M(t_i)/N$. Dividing by $N_0$ does not
remove survival. Dividing instead by the surviving number $N_0p(t_i)$ would
also amplify the noise by $1/p(t_i)$ and defines a different record.

\subsection*{Measurement strength needs a Sr-specific calibration}
For a fixed isolated transition and geometry, the expected scaling is
\begin{equation}
\kappa=q_{\mathrm{meas}}\frac{\sigma_{\mathrm{ref}}}{A}\gamma_0,
\qquad
\OD_0=\frac{N_0\sigma_{\mathrm{ref}}}{A},
\qquad
\operatorname{Var}(m_i)=
\frac{1}{q_{\mathrm{meas}}N_0\OD_0\gamma_0\dt}.
\label{eq:noise-scaling}
\end{equation}
Here $A$ is the beam mode area, $\sigma_{\mathrm{ref}}$ is a consistently
defined reference cross section, and $q_{\mathrm{meas}}>0$ contains the vector
polarizability, optical normalization, and detection efficiency. Its numerical
value has not been derived in this discussion. Measurement strength involves
the \emph{square} of the Faraday signal coupling; the note's prefactor written
as $C^{(1)}$ needs its convention specified.

The weak return branch also sets the dipole strength of the probed transition.
A closed-two-level resonant cross section cannot be substituted without
accounting for the branching and angular factors. The reference OD in
Eq.~\eqref{eq:noise-scaling} should not be confused with an off-resonant
absorption optical depth, which itself depends on detuning and state.

For an initially uncorrelated spin coherent ensemble, an indicative
weak-backaction condition is
$\kappa T N_0 J/2\ll1$. Under Eq.~\eqref{eq:noise-scaling}, this becomes
$q_{\mathrm{meas}}\OD_0\gamma_0 T J/2\ll1$, up to the evolving spin variance.
The note's $\OD\gamma T\ll1$ is the corresponding scaling estimate.

\newpage
\section{Detuning, resources, and the estimation objective}

The proposed additional control is the detuning dependence of
Eq.~\eqref{eq:beta}. Writing probe intensity as $I_L$, the isolated-line
scalings are
\begin{equation}
|\beta|\propto\frac{I_L}{|\Delta_L|},\qquad
\gamma_0\propto\frac{I_L}{\Delta_L^2},\qquad
\kappa\propto\gamma_0.
\end{equation}
Consequently, the effect of increasing $|\Delta_L|$ depends on which resource
is held fixed:
\begin{center}
\begin{tabularx}{\textwidth}{@{}lYY@{}}
\toprule
Fixed quantity & Required intensity & Effect of larger detuning\\
\midrule
$I_L$ & Constant & Weaker tensor dynamics, less scattering, weaker readout.\\
$|\beta|$ & $I_L\propto|\Delta_L|$ & Less scattering, but also weaker readout.\\
$\gamma_0$ & $I_L\propto\Delta_L^2$ & Stronger tensor dynamics at fixed reference scattering and readout strength.\\
\bottomrule
\end{tabularx}
\end{center}

The last comparison directly tests whether extra nonlinear dynamics improves
vector estimation at a fixed measurement-noise scale. With state-dependent
loss, fixing $\gamma_0$ fixes the loss \emph{operator scale}; the actual
survival trajectory can still change as the Hamiltonian changes. Increasing
the tensor strength does not by itself prove a Fisher-information gain.

These scalings require weak excitation and a sufficiently isolated transition,
schematically
\begin{equation}
\Gamma,\ |\Omega_L|\ll|\Delta_L|
\ll\text{separation to relevant neighboring optical lines}.
\end{equation}
At larger detunings, contributions from other excited manifolds must be
included. Available laser power and acceptable heating also constrain an
experimental detuning sweep. Arbitrarily large detuning is not a consequence
of this reduced model.

\subsection*{Accumulated information and useful measurement duration}
For a fixed control and probe setting, let
$\mu_i(\bm\omega)=\Tr[\Jk\rho(t_i;\bm\omega)]$ and let the known shot-noise
variance be $\sigma_i^2$, independent of the unknown field. Then
\begin{equation}
F(T)=\sum_{t_i\le T}\frac{1}{\sigma_i^2}
\bigl(\nabla_{\bm\omega}\mu_i\bigr)
\bigl(\nabla_{\bm\omega}\mu_i\bigr)^{\mathsf T}.
\label{eq:fisher}
\end{equation}
Each additional measurement contributes a positive semidefinite matrix.
Accumulated information cannot decrease when a record is extended under
otherwise identical conditions, although signal loss may make the gain very
small. Where $F$ is invertible, the associated Cramer--Rao bound cannot worsen.

A finite optimal duration can arise when optimizing information per total
experimental time, including repeated preparation and dead time, or when
jointly changing the probe settings at fixed duration. Thus
$\gamma_{\mathrm{loss}}T\sim1$ is a useful heuristic timescale for uniform
loss, not a general maximum of accumulated Fisher information. For vector
estimation, a scalar objective such as $\Tr F^{-1}$ must also be specified.

\newpage
\section{Comparison with v7 and proposed next decisions}

The current script is
\path{multiparameter_trotter_v7_loss_least_squares_over_time.py}.
Its loss generator already contains a term $-\{K_L,\rho\}/2$, with
$K_L=\gamma_s(a_L\id+c_LJ_z^2)$, alongside within-manifold spin-flip channels.
Thus the trace-decreasing formalism can be reused; the atomic coefficients,
probe parametrization, and any return channels require a new model.

\begin{center}
\begin{tabularx}{\textwidth}{@{}lYY@{}}
\toprule
Ingredient & Current v7 & Proposed Sr model\\
\midrule
Sensing manifold & Cs $F=3$ or $4$ & Sr $J=2$\\
Irreversible dynamics & Spin flips and Raman loss & Dominant loss; optional $1.88\%$ return branch\\
Tensor/scattering ratio & Fixed atomic preset & Detuning-dependent ratio\\
Probe parameters & Sweep positive $\beta$ magnitude & Intensity and detuning set shift, scattering, and noise\\
Loss geometry & Preset $a_L+c_LJ_z^2$ & $Q=(4\id-\Je^2)/10$\\
Readout & Unnormalized signal plus shot noise & Same signal convention; Sr-specific optical calibration needed\\
\bottomrule
\end{tabularx}
\end{center}

\textbf{Recommended starting point, not yet an agreed choice.}
Use Eqs.~\eqref{eq:hamiltonian}, \eqref{eq:scales}, \eqref{eq:beta},
\eqref{eq:Q}, \eqref{eq:loss}, and \eqref{eq:record}. Retain state-dependent
loss while initially treating all scattering as loss. Compare it with a
uniform-loss model to determine whether state-selective survival changes the
inferred metrological advantage. The five-dimensional spin space makes this
comparison modest in size.

The remaining decisions are:
\begin{itemize}
\item Retain state-dependent loss immediately, or start with a scalar
approximation? If scalar, specify how its rate is calibrated: the isotropic
average is $\Tr Q/5=1/5$, while an initial spin coherent state transverse to
$\boldsymbol\epsilon$ has $\langle Q\rangle=3/10$.
\item Choose polarization, propagation, initial spin, and control axes together.
\item Derive the Faraday coefficient and cross-section normalization before
assigning a numerical OD or shot-noise level.
\item Specify the detuning/power constraint and the Fisher-information objective
for the first comparison. Add the small return branch only when its effect is
needed at the desired accuracy.
\end{itemize}

\begin{thebibliography}{9}
\small
\bibitem{note}
IHD, handwritten note, \emph{Vector Magnetometry}, September 2026,
pp.~1--3 (page~4 is blank).
Repository file: \texttt{Vector Magnetometry - 20260926\_151513.pdf}.

\bibitem{masson}
S.~J.~Masson, J.~P.~Covey, S.~Will, and A.~Asenjo-Garcia,
\emph{Dicke Superradiance in Ordered Arrays of Multilevel Atoms},
PRX Quantum \textbf{5}, 010344 (2024), Fig.~1(c) and Table~II.
\href{https://doi.org/10.1103/PRXQuantum.5.010344}{doi:10.1103/PRXQuantum.5.010344}.
Repository file: \path{PRXQuantum.5.010344.pdf}.

\bibitem{deutsch}
I.~H.~Deutsch and P.~S.~Jessen,
\emph{Quantum control and measurement of atomic spins in polarization
spectroscopy}, Optics Communications \textbf{283}, 681--694 (2010).
See Eq.~(14) for the tensor convention and Sec.~3.2.1 for measurement strength.
\href{https://doi.org/10.1016/j.optcom.2009.10.059}{doi:10.1016/j.optcom.2009.10.059}.
\end{thebibliography}

\end{document}
